\scanpage{236}

\[
\begin{array}{ll}
(29)&\left\{
\begin{aligned}
r_\infty&=\tau_\infty=\begin{pmatrix}-1&0\\0&1\end{pmatrix}\\
r_0&=\tau'_\infty=\begin{pmatrix}0&1\\1&0\end{pmatrix}\\
r_1&=-\tau'_0=\begin{pmatrix}-1&1\\0&1\end{pmatrix}
\end{aligned}
\right.
\end{array}
\]
\margin{$\left[\begin{array}{ll}
(30)&\left\{
\begin{aligned}
r_0(\rho)&=\rho^{-1},\quad r_0(\sigma)=\sigma^{-1}\\
r_1&=\rho^{-1}r_0=r_0\rho\\
r_\infty&=\sigma r_0=r_0\sigma^{-1}
\end{aligned}\right.
\end{array}\right.$}

Note that $(\widetilde\rho,\widetilde\tau')$ \unsure{generate} a group isomorphic to $D_\infty$, which I denote by $D_{\widetilde\rho}$,
likewise $(\widetilde\sigma,\widetilde\tau')$ \unsure{generates|generate} a group $D_{\widetilde\sigma}$,
isomorphic to $D_\infty$. Finally $(\omega_0,\widetilde\tau')$ \unsure{generate} a group $D_{\omega_0}$, again isomorphic to $D_\infty$. We have a diagram
\diagram{
$(31)\quad$
\begin{tikzcd}[ampersand replacement=\&,column sep=large,row sep=large]
\&D_{\omega_0}\arrow[dl,"i"']\arrow[dr,"j"]\&\\
D_{\widetilde\rho}\&\&D_{\unsure{\widetilde\sigma|\widetilde\rho}}
\end{tikzcd}
\qquad
$\begin{aligned}
&\left\{\begin{aligned}
i(\omega_0)&=\widetilde\rho^{-3}\\
i(\widetilde\tau')&=\widetilde\tau'
\end{aligned}\right.\\[4pt]
&\left\{\begin{aligned}
j(\omega_0)&=\widetilde\sigma^2\\
j(\unsure{\widetilde\tau'})&=\widetilde\tau'
\end{aligned}\right.
\end{aligned}$
}
and corollary 1 is written in the form
\[
\begin{array}{ll}
(32)&\mathrm{Gl}(2,\mathbb Z)^{\sim}\simeq D_{\widetilde\rho}*_{D_{\omega_0}}D_{\widetilde\sigma}
\end{array}
\]
(amalgamated \unsure{sum} of groups).

Now we can take in $D_{\widetilde\rho},D_{\widetilde\sigma}$ the involutive generators
\[
\begin{array}{ll}
(33)&\left\{\begin{array}{lll}
\widetilde r_0=\widetilde\tau'=\widetilde\tau'_\infty&\text{and}&
\widetilde r_1\overset{\mathrm{dfn}}{=}\widetilde\rho^{-1}\widetilde r_0=\widetilde r_0\widetilde\rho\quad\text{in }D_{\widetilde\rho}\\
\widetilde r_0=\cdots\cdots&\text{and}&
\widetilde r_\infty\overset{\mathrm{dfn}}{=}\unsure{\widetilde\sigma\,\widetilde r_0}=\widetilde r_0\widetilde\sigma^{-1}
\end{array}\right.
\end{array}
\]

\margin{235}

\begin{center}
235\\
\texttt{https://grothendieck.umontpellier.fr}\\
Université de Montpellier, UMR CNRS 5149, Case courrier 051, Place Eugène Bataillon\\
34095 Montpellier cedex 5 - France
\end{center}

\scanpage{237}
\mspage{766}

and we then find

\uline{Corollary 2} $\mathrm{Gl}(2,\mathbb Z)^{\sim}$ admits a presentation (with three involutive generators $\widetilde r_0,\widetilde r_1,\widetilde r_\infty$):
\[
\begin{array}{ll}
(34)&\boxed{\widetilde r_0^2=\widetilde r_1^2=\widetilde r_\infty^2=1,\quad
\Bigl(\underbrace{\widetilde r_0^{\deleted{\unsure{-1}}}\widetilde r_1}_{\widetilde\rho}\Bigr)^3
\overset{\deleted{\unsure{=}}}{\cdot}
\Bigl(\underbrace{\widetilde r_\infty\widetilde r_0^{\deleted{\unsure{-1}}}}_{\widetilde\sigma}\Bigr)^2=1}\\[-2pt]
&\hspace{13em}\deleted{\underbrace{\hspace{9em}}_{\unsure{\omega_0}}}
\end{array}
\]

\margin{236}

\begin{center}
236\\
\texttt{https://grothendieck.umontpellier.fr}\\
Université de Montpellier, UMR CNRS 5149, Case courrier 051, Place Eugène Bataillon\\
34095 Montpellier cedex 5 - France
\end{center}

\scanpage{238}

$4^\circ)$ \uline{The isomorphism} $\mathrm{Gl}(2,\mathbb Z)^{\sim}\simeq S\mathcal T_{1,1}$

\margin{$\unsure{\checkmark}$}
Put
\[
\begin{array}{ll}
(35)&\deleted{\unsure{X_1}}\quad X_1=\mathbb R^2/\mathbb Z^2\simeq(\mathbb R/\mathbb Z)^2\simeq\mathbb U\times\mathbb U.
\end{array}
\]
(model surface of genus $1$), write its composition law additively, and let
\[
\begin{array}{ll}
(36)&a_{1,0}=0\in X_1,\quad X_1\setminus\{a_{1,0}\}=X_1\setminus\{0\}=U_{1,1},
\end{array}
\]
model surface of type $(g,\nu)=(1,1)$.

The group $\mathrm{Gl}(2,\mathbb Z)$ is identified with the group of automorphisms of the topological group $X_1$, and thus operates on $X_1$ leaving fixed the origin $\deleted{a_{1,0}}\ 0$, hence it operates on $U_{1,0}$, whence a homomorphism
\diagram{
$(37)\quad$
\begin{tikzcd}[ampersand replacement=\&,column sep=large,row sep=tiny]
\mathrm{Gl}(2,\mathbb Z)\arrow[r,"\sim"]\&\mathcal T_{1,1}=\mathcal T(\unsure{U_{1,1}})\\
\&\begin{gathered}\text{group of isotopy classes}\\
\text{of homeomorphisms}\\
\text{of }U_{1,1},\end{gathered}
\end{tikzcd}
}
which we know to be an isomorphism. Moreover, $G=\mathrm{Gl}(2,\mathbb Z)$ operates \unsure{differentiably|differentially} on $\deleted{\unsure{\mathcal T|\tau}}\ X_1$, and the action on the \uline{tangent} space is identified with the tautological action of $\mathrm{Gl}(2,\mathbb Z)\subset\mathrm{Gl}(2,\mathbb R)$ on $\mathbb R^2$. Thus the group denoted by $\widetilde G$ in

\margin{237}

\begin{center}
237\\
\texttt{https://grothendieck.umontpellier.fr}\\
Université de Montpellier, UMR CNRS 5149, Case courrier 051, Place Eugène Bataillon\\
34095 Montpellier cedex 5 - France
\end{center}

\scanpage{239}
\mspage{767}

\unsure{$1^\circ$)} \unsure{e} is canonically isomorphic, as an extension of $G$ by $T_0=\mathbb Z$ (N.B. $\mathbb R^2$ and its quotient $X_1$ are canonically oriented) to the group $\mathrm{Gl}(2,\mathbb Z)^\sim$ \unsure{of} $3^\circ$.

Thus we find \deleted{the \unsure{Theorem}} \inserted{(cf. proposition p. 761)} the

\uline{Theorem} The inverse image under the isomorphism (37) of the extension $S\deleted{\unsure{*}}\mathcal T_{1,1}$ of $\mathcal T_{1,1}$ by $\mathbb Z$ is canonically identified with $\mathrm{Gl}(2,\mathbb Z)^\sim$, so it admits the \inserted{two} explicit presentations \unsure{of} $3^\circ$
\[
\begin{aligned}
\text{a)}\quad &
\underset{\deleted{\unsure{\omega\widetilde\sigma}}}{\widetilde\rho^{3}\cdot\widetilde\sigma^{2}}
=\underset{\deleted{\unsure{\omega_0}}}{1},\quad
\unsure{\widetilde\tau'}(\widetilde\rho)=\widetilde\rho^{-1},\quad
\widetilde\tau'(\widetilde\sigma)=\widetilde\sigma^{-1},\quad
\widetilde\tau'^{\,2}=1
\\[1ex]
\text{b)}\quad &
\widetilde r_0^{2}=\widetilde r_1^{2}=\widetilde r_\infty^{2}=1,\quad
\underset{\deleted{\unsure{\omega_0}}}{
(\widetilde r_0^{\deleted{\unsure{-1}}}\widetilde r_1)^3\cdot
(\widetilde r_\infty\widetilde r_0)^2}=1.
\end{aligned}
\]

Let us now interpret $\mathcal T_{1,1}$ and $S\mathcal T_{1,1}$ in terms of the fundamental group of $U_{1,1}$, which we \unsure{shall compute} by choosing (judiciously ?) a base point $a_{1,1}$ in $U_{1,1}$,
\[
\text{(38)}\qquad
\Pi_{1,1}\overset{\mathrm{dfn}}=\pi_1(U_{1,1},a_{1,1}),
\]
\[
\text{(39)}\qquad
\mathcal T_{1,1}\simeq\mathrm{Aut\,ext}_{\mathrm{lac}}(\Pi_{1,1})
\]
\margin{238}

\begin{center}
238\\
\texttt{https://grothendieck.umontpellier.fr}\\
Université de Montpellier, UMR CNRS 5149, Case courrier 051, Place Eugène Bataillon\\
34095 Montpellier cedex 5 - France
\end{center}

\scanpage{240}

\margin{\unsure{I use} the \unsure{notation} $\ell'_0$ instead of $\ell_0$, since $\ell_0$ already denotes a well-determined element of $\mathrm{Gl}(2,\mathbb Z)$ !}

We have a presentation of $\Pi_{1,1}$ by
\[
\text{(40)}\qquad
\Pi_{1,1}=\{x,y,\ell'_0\mid [x,y]\ell'_0=1\}
\]
and denoting by $L$ the subgroup\inserted{\unsure{-loops}} of $\Pi_{1,1}$ generated by $\ell'_0$, we shall have
\[
\text{(41)}\qquad
S\mathcal T_{1,1}\simeq
\mathrm{Norm}_{\mathrm{Aut}_{\mathrm{lac}}(\Pi_{1,1})}(L).
\]

Thus we should make explicit the action
\deleted{\unsure{of} $\unsure{S}$}
\diagram{\begin{tikzcd}[ampersand replacement=\&,column sep=large]
S\mathcal T_{1,1}\simeq\mathrm{Gl}(2,\mathbb Z)^\sim
\arrow[r] \& \mathrm{Aut}_{\mathrm{lac}}(\Pi_{1,1}),
\end{tikzcd}}
by the action of the generators $\widetilde\rho,\widetilde\sigma,\widetilde\tau'$ or $\widetilde r_0=\widetilde\tau',\widetilde r_1,\widetilde r_\infty$ --- action which must \unsure{normalize} $L$, and more precisely, transform $\ell'_0$ into $\ell'_0$ (for $\widetilde\rho,\widetilde\sigma$) or into $\ell_0'^{-1}$ (for $\widetilde\tau'=\widetilde r_0,\widetilde r_1,\widetilde r_\infty$). Thus we must make explicit the action of the generators on $x$ and on $y$. (N.B. $\Pi_{1,1}$ is the free group having as generators $x$ and $y$ \unsure{!})

Let us consider the principal covering $U_{1,1}^{\prime\,\unsure{a}}$ with group $\mathbb Z^2$ of $U_{1,1}=X_1-\{0\}$, induced by the principal \unsure{covering} $\mathbb R^2$ of $X_1=\mathbb R^2/\mathbb Z^2$, --- the inverse image of $U_{1,1}$ is $\mathbb R^2\setminus\mathbb Z^2$. Let us choose a base point \deleted{$\unsure{a}$}
\margin{239}

\begin{center}
239\\
\texttt{https://grothendieck.umontpellier.fr}\\
Université de Montpellier, UMR CNRS 5149, Case courrier 051, Place Eugène Bataillon\\
34095 Montpellier cedex 5 - France
\end{center}